\newcolumntype{a}{>{\columncolor{RoyalBlue!50}}l}
\newcolumntype{b}{>{\columncolor{RoyalBlue!50}}p{1.2cm}}
\newcolumntype{d}{>{\columncolor{RoyalBlue!50}}p{.6cm}}
\newcolumntype{R}[1]{%
    >{\columncolor{RoyalBlue!50}\adjustbox{angle=#1}\bgroup}%
    l<{\egroup}|%
}
\newcommand*\rotCol{\multicolumn{1}{R{90}}}

\lstset{style=sql, basicstyle=\ttfamily\scriptsize}

\title[Datenbanken] {Datenbanken}

\subtitle{\emoji{floppy-disk} \textbf{\gls{SQL}: mehrere Tabellen}}

\institute[KZN] % (optional)
{
  \inst{1}%
  Fachschaft Informatik\\
  Kantonsschule Zürich-Nord
  \and
  \inst{2}%
  Fachschaft Informatik\\
  Kantonsschule im Lee
}

\date[\today] % (optional)

\logo{\includegraphics[width=.8cm]{Logos/FreeFlower.pdf}}

\titlegraphic{\includegraphics[width=\textwidth]{"Grundlagen_Info/06_Datenbanken/Figures/title-emojis"}}

\frame{\titlepage}


\begin{frame}[fragile]{Weshalb braucht es mehrere Tabellen?}{Analogie: Wörterbuch}
Angenommen, Sie lesen einen Text in deutscher Sprache. Wäre es praktisch, wenn jedes Wort in einer Klammer erklärt wäre?\\[.5cm]


\begin{tcolorboxTemplate}[title=Beispiel]
\tiny\textbf{Ich} (\enquote{das Selbst, dessen man sich bewusst ist und mit dem man sich von der Umwelt unterscheidet}) \textbf{ging} (gehen, \enquote{schrittweises Sichfortbewegen auf den Füssen in aufrechter Haltung}) \textbf{nach} (\enquote{drückt aus, dass jemand jemandem, einer Sache folgt, nachgeht}) \textbf{draussen} (\enquote{ausserhalb eines Raumes, Gebäudes})...
\end{tcolorboxTemplate}


\normalsize
Alternative: Nachschlagewerk / Wörterbuch!

\end{frame}

\begin{frame}[fragile]{Weshalb braucht es mehrere Tabellen?}{Analogie: Soziale Medien}
\begin{center}
Wäre es praktisch, wenn Sie bei Ihrer Tiktok-Freundesliste immer auch gleich all deren Posts, Kommentare und alles Weitere sehen könnten?
\end{center}
\end{frame}

\begin{frame}{Tabellen-Organisation}
\Wider{\begin{figure}[H]
\centering
\includegraphics[width=\linewidth]{"Grundlagen_Info/06_Datenbanken/Figures/schema MM"}
\end{figure}}
\begin{itemize}
\item Was bedeuten die blauen Pfeile und die Schlüssel?
\end{itemize}
\end{frame}


\begin{frame}[fragile]{Primärschlüssel einer Tabelle}

\begin{table}
\centering
\begin{tabular}{a|l|l}
\emoji{key} \underline{mID} & Name   & ... \\\midrule
1  & Müller  & ...  \\
2  & Schmidt & ...  \\
3  & Kaufmann  & ...\\
... & ... & ...
\end{tabular}
\caption*{Tabelle \enquote{\textbf{Angestellte}}}
\end{table}
\begin{center}
\textbf{Frage}: Weshalb hat (fast) jede Tabelle einen Primärschlüssel \emoji{key}?
\end{center}
\end{frame}

\begin{frame}[fragile]{Fremdschlüssel einer Tabelle}

\begin{columns}
\begin{column}{.5\textwidth}
\begin{table}[H]
\begin{tabular}{l|l|a}
\emoji{key} \underline{mID} & Name   & \tikzmarknode{fkey}{\underline{aID} \emoji{up-right-arrow}} \\\midrule
1  & Müller  & \tikzmarknode{31o}{31}  \\
2  & Schmidt & \tikzmarknode{32o1}{32}  \\
3  & Kaufmann  & \tikzmarknode{32o2}{32}\\
... & ... & ...
\end{tabular}
\caption*{Tabelle \enquote{\textbf{Angestellte}}}
\end{table}
\end{column}
\begin{column}{.5\textwidth}
\begin{table}[H]
\begin{tabular}{a|l}
\tikzmarknode{pkey}{\emoji{key} \underline{aID}} & Abteilung \\\midrule
\tikzmarknode{31t}{31}    & Verkauf   \\
\tikzmarknode{32t}{32}    & Technik   \\
\tikzmarknode{33t}{33}    & Marketing \\
... & ...
\end{tabular}
\caption*{Tabelle \enquote{\textbf{Abteilungen}}}
\end{table}
\end{column}
\end{columns}


\only<2->{
\begin{tikzpicture}[overlay,remember picture]
% title
 \draw[red,-stealth] (fkey) to[bend left] (pkey);
% values
  \draw[red,-stealth] (31o) -- (31t);
  \draw[red,-stealth] (32o1) -- (32t);
  \draw[red,-stealth] (32o2) -- (32t);
\end{tikzpicture}}

\begin{center}
\only<1>{
\textbf{Frage}: Was könnte ein Fremdschlüssel \emoji{up-right-arrow} sein? \\

Weshalb könnten Fremdschlüssel benötigt werden?
}
\only<2>{
\textbf{Frage}: Was ist die Abteilung von Frau Kaufmann?
}
\end{center}

\end{frame}

\begin{frame}{Beziehungen darstellen: \glsentryshort{ERM}}{Chen-Notation}
  \begin{center}
    \begin{tikzpicture}[auto]
      \node[entity] (user) at (0,0) {users};
      \node[relationship] (owns) at (3,0) {owns};
      \node[entity] (photo) at (6,0) {photos};
      \draw (user) -- (owns) node[midway, above] {1};
      \draw (owns) -- (photo) node[midway, above] {n};
    \end{tikzpicture}
  \end{center}
  \begin{itemize}
    \item Rechteck: Tabelle (\textit{Entity}), Raute: Beziehung (\textit{Relationship})
    \item Kardinalitäten: 1:1, 1:n (ein Mitglied, viele Fotos), n:m (Likes: viele Mitglieder, viele Fotos)
    \item \enquote{0:1} heisst \enquote{keine oder eine}
  \end{itemize}
\end{frame}

\begin{frame}{Auftrag}{\glsentryshort{ERM}}
  \begin{itemize}
    \item \readbox{Skript, Abschnitt 3.1}
    \item \aufgabebox{Aufgaben 3.1, 3.2 im Skript}
  \end{itemize}
\end{frame}

\begin{frame}[fragile]{Mehrere Tabellen miteinander verbinden}{Mit \lstinline[style=sql]|JOIN| und \lstinline[style=sql]|ON| (falls die Kolonnennamen ungleich sind)}
	\begin{lstlisting}
		SELECT *
		FROM Angestellte JOIN Abteilungen
        ON Angestellte.aID = Abteilungen.ID
	\end{lstlisting}
	\vspace{-.3cm}
	\scriptsize
	\begin{columns}
		\begin{column}{.35\textwidth}
			\begin{table}[H]
				\begin{tabular}{l|l|a}
					\emoji{key} \underline{mID} & Name   & \underline{aID} \emoji{up-right-arrow} \\\midrule
					1  & Müller  & 31  \\
					2  & Schmidt & 32  \\
					3  & Kaufmann  & 32\\
					... & ... & ...
				\end{tabular}
				\caption*{Tabelle \enquote{\textbf{Angestellte}}}
			\end{table}
			\vspace{-.5cm}\begin{table}[H]
				\begin{tabular}{a|l}
					\emoji{key} \underline{ID} & Abteilung \\\midrule
					31    & Verkauf   \\
					32    & Technik   \\
					33    & Marketing \\
					... & ...
				\end{tabular}
				\caption*{Tabelle \enquote{\textbf{Abteilungen}}}
			\end{table}

		\end{column}
		\begin{column}{.1\textwidth}
			\faArrowRight
		\end{column}
		\begin{column}{.55\textwidth}
			\begin{table}[H]
				\begin{tabular}{l|l|a|a|l}
					mID & Name   & \rotCol{Angestellte.aID} & \rotCol{Abteilungen.ID} & Abteilung \\\midrule
					1  & Müller  & 31  &31  & Verkauf\\
					2  & Schmidt & 32 &32 & Technik\\
					3  & Kaufmann  & 32 &32 & Technik \\
					... & ... & ... & ...
				\end{tabular}
			\end{table}
		\end{column}
	\end{columns}
\end{frame}

\begin{frame}[fragile]{Verteilte Tabellen}


\begin{table}[H]
\centering
\tiny
\adjustbox{max width=\textwidth, max totalheight=.62\textheight}{
\begin{tabular}{l|p{7cm}}
\textbf{Tabelle} & \textbf{Spalten} \\\midrule\midrule
Lieferant&	\emoji{key} LieferantenNr, Firma, Kontaktperson, Position, Strasse, Ort, Region, PLZ, Land, Telefon, Telefax, Homepage\\\midrule
Kunde &	\emoji{key} KundenCode, Firma, Kontaktperson, Position, Strasse, Ort, Region, PLZ, Land, Telefon, Telefax\\\midrule
Versandfirma &	\emoji{key} FirmenNr, Firma, Telefon\\\midrule
Personal &	\emoji{key} PersonalNr, Nachname, Vorname, Position, Anrede, Geburtsdatum, Einstellung, Strasse, Ort, Region, PLZ, Land, TelefonPrivat, DurchwahlBüro, Bemerkungen, Vorgesetzter\\\midrule
Kategorie 	&\emoji{key} KategorieNr, Kategoriename, Beschreibung\\\midrule
Artikel 	&\emoji{key} ArtikelNr, Artikelname, LieferantenNr \emoji{up-right-arrow}, KategorieNr \emoji{up-right-arrow}, Liefereinheit, Einzelpreis, Lagerbestand, BestellteEinheiten, Mindestbestand, Auslaufartikel\\\midrule
Bestellung 	& \emoji{key} BestellNr, KundenCode \emoji{up-right-arrow}, PersonalNr \emoji{up-right-arrow}, Bestelldatum, Lieferdatum, Versanddatum, FirmenNr \emoji{up-right-arrow}, Frachtkosten, Empfänger, Strasse, Ort, Region, PLZ Bestimmungsland\\\midrule
Bestelldetails &	BestellNr \emoji{up-right-arrow}, ArtikelNr \emoji{up-right-arrow}, Einzelpreis, Anzahl, Rabatt
\end{tabular}}
\end{table}
\vspace{-.3cm}
\begin{figure}[H]
\centering
\includegraphics[height=.8cm]{Grundlagen_Info/06_Datenbanken/Figures/Migros}
\includegraphics[height=.8cm]{Grundlagen_Info/06_Datenbanken/Figures/Coop}
\end{figure}

\end{frame}


\begin{frame}{Auftrag}{Verbinden von Tabellen}
\textbf{Moodle}
\begin{itemize}
\item \aufgabebox{Moodle 3.1--3.3}
\item \challengebox{Moodle 3.4}
\item \challengebox{Abgabe auf Moodle (InstaHub, \enquote{\texttt{JOIN}})}
\end{itemize}
\textbf{Skript (Abschnitte 3.2 und 3.3)}
\begin{itemize}
\item \aufgabebox{Aufgaben 3.3--3.5, 3.7, 3.8, 3.10}
\item \challengebox{Aufgaben 3.6, 3.9}
\end{itemize}
\end{frame}
